\section{Notation, intermediate checks, and reading exercises}
\label{app:reader}

\subsection{A notation sheet}

\begin{longtable}{p{25mm}p{112mm}}
\toprule
Symbol&Meaning\\
\midrule
\endhead
\(k\)&Fixed target exponent in \(\pi^k\).\\
\(L,t_j\)&Denominator of the interval shape and its distinct offsets.\\
\(\nu_j\)&Multiplicity of interval type \(j\).\\
\(M,m\)&Total multiplicity \(M=\sum_j\nu_j\); logarithmic degree \(m=M-1\).\\
\(n=LN_0\)&Unbounded approximation index; \(N_0\) is a positive integer.\\
\(c_j\)&Relative interval length \(1-2t_j/L\).\\
\(Q_n,C_n\)&Rational-kernel denominator and the integer normalization factor.\\
\(d_n\)&Least common multiple of \(1,\ldots,n\).\\
\(\Pi_n,\gamma\)&Available common divisor and its exponential rate, or a certified lower rate when specified.\\
\(\lambda,\Theta\)&\(\lambda=\log\ii=\ii\pi/2\), \(\Theta=\pi/2\).\\
\(Y\)&Positive contour height after rescaling by \(n\); not the polynomial indeterminate in a particular formula.\\
\(D_0,G_0\)&Decay and growth rates before division by the common divisor.\\
\(D,G\)&Usable rates after division, possibly weakened to conservative rational bounds.\\
\(\Delta_k\)&Effective gap \(D-(k-1)G\), required to be positive.\\
\(s,v_j\)&Common scale and fixed base weights, so \(\nu_j=sv_j\).\\
\(V,S\)&Base multiplicity \(\sum_jv_j\) and weighted base length \(\sum_jv_jc_j\).\\
\(p,q\)&Arbitrary rational-approximation numerator and positive denominator.\\
\(\theta,y\)&Nearby algebraic numbers \((p/q)^{1/k}\) and \(\ii\theta/2\).\\
\bottomrule
\end{longtable}

The notation \(Y\) is conventional both for a contour height and a
polynomial variable. In a computation it is safer to use different
names. Similarly, the small finite field-independent degree \(k\)
of a defining binomial must not be confused with the possibly
large auxiliary degree \(m\).

\subsection{Deriving the exponent without suppressing the index choice}

Let \(\eta>0\) be the common slack in the three asymptotic
estimates, and set \(\Delta_\eta=D-(k-1)G-k\eta>0\).
The resultant lower bound has the form
\[
|U(y)|\ge(2^kq)^{-m}e^{-(k-1)(G+\eta)n}.
\]
Choose a multiple of \(L\) so that
\(n\ge(m\log(2^kq)+\log2)/\Delta_\eta\). The remainder
\(e^{-(D-\eta)n}\) is then at most half this lower bound.
The divided difference contributes one more growth factor,
giving
\[
|\lambda-y|\ge\tfrac12(2^kq)^{-m}e^{-k(G+\eta)n}.
\]
Since rounding increases \(n\) by less than \(L\), its effect is
a positive constant independent of \(q\). The power of \(q\) is
\[
m+\frac{mk(G+\eta)}{\Delta_\eta}
=\frac{m(D+G)}{\Delta_\eta}.
\]
This calculation explains both the denominator and the numerator
of the final bound. It also shows why omitting the divided
difference, or forgetting a conjugate in the resultant, changes
the exponent rather than merely the prefactor.

\subsection{Short exercises with solutions or hints}

\paragraph{1. Rational rescaling.}
Show that \(\mu(a\xi/b)=\mu(\xi)\) for nonzero integers
\(a,b\). A rational approximation to one side gives one to the
other after multiplying numerator and denominator by fixed
integers. After reduction, the denominator changes by a factor
bounded above and below independently of the approximation.
Apply this to \(\beta(3)=\pi^3/32\) and
\(\zeta(4)=\pi^4/90\). Notice that taking powers does not have
this bounded-denominator property in both directions.

\paragraph{2. The elementary arctangent inequality.}
Differentiate \(\arctan u-u+u^3/3\). Its derivative is
\(u^4/(1+u^2)\ge0\), and the value at zero is zero.
This proves the inequality used in the root-free estimate without
any alternating-series restriction such as \(u\le1\).

\paragraph{3. Why the exact saddle is not necessary.}
For fixed multiplicities, differentiate
\(\Phi(Y)=\sum_j\nu_jF_{c_j}(Y)-\Theta Y-m\).
The second derivative is negative. Thus the unique critical
point is a global maximum, but every \(Y>1/2\) still provides
a valid lower decay rate. A rational trial height trades a small
amount of sharpness for shorter numerical certificates.

\paragraph{4. The degree obstruction.}
Rewrite \(E_k=m(D+G)/(D-(k-1)G)\) as
\(m(1+kG/\Delta_k)\). If \(G\ge0\), then \(E_k\ge m\).
In particular, no refinement of these same rates at fixed
\(m=171\) can yield an exponent below 100 through this formula.
A new transfer argument or a smaller auxiliary degree would be
needed. This does not rule out a different construction.

\paragraph{5. A smaller true algebraic degree.}
Take \(p/q=(a/b)^k\) with \(a,b>0\). Then
\(y=\ii a/(2b)\) already belongs to \(\QQ(\ii)\).
Its minimal polynomial over that field has degree one, even
though the defining binomial has degree \(k\). Explain why the
minimal-polynomial resultant remains nonzero for a selected row,
and why a product over all roots of the binomial need not be.

\paragraph{6. A common divisor and the gap.}
Keeping \(D_0,G_0\) fixed, increase the usable divisor rate by
\(h>0\). The new pair is \((D+h,G-h)\), so the gap increases
by \(kh\) while \(D+G\) is unchanged. This is the precise
reason the same divisor has to be used in both the analytic
remainder and the coefficient height.

\paragraph{7. What a finite numerical test can establish.}
Suppose \(-n^{-1}\log|Q_n\beta(3)-P_n|\) is positive at
each of the first one hundred indices. This does not prove that
the integer linear forms tend to zero, that their asymptotic
rate is positive, or that the nonproportionality condition needed
for every denominator holds. State the additional asymptotic and
determinant hypotheses before extracting an irrationality bound.
The failed experiments in Appendix~\ref{app:experiments} should
be read with exactly this distinction in mind.
